\documentclass[12pt]{article}
\usepackage{fullpage,amsmath,amssymb,amsthm,hyperref,amsfonts,enumitem}

\newtheorem{theorem}{Theorem}
\newtheorem{lemma}[theorem]{Lemma}
\newtheorem{proposition}[theorem]{Proposition}
\newtheorem{corollary}[theorem]{Corollary}
\theoremstyle{definition}
\newtheorem{definition}[theorem]{Definition}
\newtheorem{remark}[theorem]{Remark}
\newtheorem{example}[theorem]{Example}

\title{A Markov Chain with Interpolation Macdonald Stationary Distribution}
\author{}
\date{}

\begin{document}
\maketitle

\begin{abstract}
Let $\lambda = (\lambda_1 > \dots > \lambda_n \geq 0)$ be a restricted partition
(distinct parts, unique part of size $0$, no part of size $1$).  We construct
an explicit nontrivial Markov chain on $S_n(\lambda)$---the set of all
rearrangements of $\lambda$---whose unique stationary distribution is
\[
  \pi(\mu) \;=\; \frac{F^*_\mu(x_1,\dots,x_n;\,q{=}1,\,t)}{P^*_\lambda(x_1,\dots,x_n;\,q{=}1,\,t)},
  \qquad \mu \in S_n(\lambda).
\]
The chain is the \emph{inhomogeneous multispecies $t$-PushTASEP} on a ring of
$n$ sites with site-dependent rates $x_1,\dots,x_n > 0$.  We prove
stationarity, irreducibility, and nontriviality in full.
\end{abstract}

\section{Background and Definitions}

\subsection{Restricted partitions and rearrangements}

Fix $n \geq 2$.  Let $\lambda = (\lambda_1 > \lambda_2 > \dots > \lambda_n \geq 0)$
be a partition with strictly decreasing parts.  We call $\lambda$
\emph{restricted} if $\lambda_n = 0$ (there is a unique part equal to $0$) and
$\lambda_i \geq 2$ for all $i \leq n-1$ (there is no part equal to $1$).

Since $\lambda_n = 0$ and $\lambda_i \geq 2$ for $i \leq n-1$:
\begin{equation}
  |\lambda| \;:=\; \sum_{i=1}^n \lambda_i \;\geq\; 2(n-1).
\label{eq:lambda_size}
\end{equation}

The \emph{species values} of $\lambda$ are the integers
$0 = \lambda_n < \lambda_{n-1} < \dots < \lambda_1$.  We regard a state
$\mu \in S_n(\lambda)$ as placing one species on each of the $n$ sites of a
directed ring (sites labeled $1, 2, \dots, n$ with $n+1 \equiv 1$).  Let
\begin{equation}
  S_n(\lambda) \;=\;
  \{\,\mu = (\mu_1,\dots,\mu_n) :
  \{\mu_1,\dots,\mu_n\} = \{\lambda_1,\dots,\lambda_n\}\text{ as multisets}\,\}.
\label{eq:statespace}
\end{equation}
Since all parts of $\lambda$ are distinct, $|S_n(\lambda)| = n!$.

\subsection{Interpolation ASEP and Macdonald polynomials}

\begin{definition}[{\cite[Def.~1.1,~Thm.~1.15]{BDW25}}]
\label{def:interpolation}
The \emph{interpolation ASEP polynomial}
$F^*_\mu(x_1,\dots,x_n;q,t) \in \mathbb{Z}_{\geq 0}[x_1,\dots,x_n,q,t]$
associated to $\mu \in S_n(\lambda)$ is the weight-generating function
for signed multiline queues with bottom row $\mu$.  It satisfies:
\begin{enumerate}[label=(\alph*)]
  \item $F^*_\mu$ has total $x$-degree equal to $|\lambda|$, and its top
        homogeneous component in $x$ is the ASEP polynomial $F_\mu(x;q,t)$.
  \item The \emph{interpolation Macdonald polynomial}
        $P^*_\lambda(x;q,t) \in \mathbb{Z}_{\geq 0}[x_1,\dots,x_n,q,t]$ is
        the unique symmetric polynomial satisfying $P^*_\lambda(x;q,t)
        = \sum_{\mu \in S_n(\lambda)} F^*_\mu(x;q,t)$.
\end{enumerate}
At $q = 1$, both $F^*_\mu(x;1,t)$ and $P^*_\lambda(x;1,t)$ have
non-negative coefficients when viewed as polynomials in $(x_1,\dots,x_n,t)$,
and $P^*_\lambda(x;1,t) > 0$ for all $x_i > 0$ and $t \in [0,1]$.
\end{definition}

We write $F^*_\mu(x;t) := F^*_\mu(x;1,t)$ and
$P^*_\lambda(x;t) := P^*_\lambda(x;1,t)$ throughout.

\begin{lemma}[Exchange identity; {\cite[Cor.~5.2]{BDW25}}]
\label{lem:exchange}
Let $\mu \in S_n(\lambda)$ and $j \in \{1,\dots,n\}$ (indices mod $n$).
If $\mu_j > \mu_{j+1}$, then
\begin{equation}
  x_j \, F^*_\mu(x;t) \;=\; x_{j+1} \, F^*_{s_j\mu}(x;t),
\label{eq:exchange}
\end{equation}
where $s_j\mu$ is the composition obtained from $\mu$ by swapping positions
$j$ and $j+1$.
\end{lemma}

\begin{remark}\label{rmk:exchange_use}
When $j \in \mathcal{A}(\mu) := \{j : \mu_j < \mu_{j+1}\}$ (an ascent of
$\mu$ at $j$), the composition $s_j\mu$ satisfies $(s_j\mu)_j =
\mu_{j+1} > \mu_j = (s_j\mu)_{j+1}$, i.e., $s_j\mu$ has a \emph{descent} at
$j$.  Applying \eqref{eq:exchange} to $s_j\mu$ at its descent $j$ gives
$x_j F^*_{s_j\mu}(x;t) = x_{j+1} F^*_\mu(x;t)$, i.e.,
\begin{equation}
  x_j \, F^*_{s_j\mu}(x;t) \;=\; x_{j+1} \, F^*_\mu(x;t)
  \qquad \text{when } j \in \mathcal{A}(\mu).
\label{eq:exchange_ascent}
\end{equation}
\end{remark}

\section{Construction of the Markov Chain}

\begin{definition}[Inhomogeneous multispecies $t$-PushTASEP]
\label{def:chain}
Let $x_1,\dots,x_n > 0$ and $0 \leq t \leq 1$.  The
\emph{inhomogeneous multispecies $t$-PushTASEP} is the continuous-time
Markov chain $(X_\tau)_{\tau \geq 0}$ on $S_n(\lambda)$ with generator $Q$:
\begin{equation}
  Q(\mu,\,\nu) \;=\;
  \begin{cases}
    x_j & \text{if } \nu = s_j\mu \text{ for some } j \in \{1,\dots,n\}
           \text{ with } \mu_j > \mu_{j+1}, \\[2pt]
    0   & \text{if } \nu \neq \mu \text{ and no such } j \text{ exists,}
  \end{cases}
\label{eq:rates}
\end{equation}
with $Q(\mu,\mu) = -\sum_{\nu \neq \mu} Q(\mu,\nu)$.  All indices are taken
mod $n$, so site $n+1 \equiv 1$.
\end{definition}

In words: at rate $x_j$, the species at site $j$ swaps with the species at
site $j+1$, provided the species at $j$ is strictly larger.  Equivalently,
a higher-species particle at site $j$ pushes clockwise past the particle at
$j+1$ at rate $x_j$.

\begin{remark}\label{rmk:welldefined}
The chain is well defined.  Since $\lambda$ has strictly distinct parts,
$\mu_j \neq \mu_{j+1}$ for all $j$ and $\mu$, so each adjacent pair is
either a strict ascent ($\mu_j < \mu_{j+1}$) or a strict descent
($\mu_j > \mu_{j+1}$).  The state space $S_n(\lambda)$ is finite, and all
off-diagonal rates are positive multiples of some $x_j > 0$.
\end{remark}

\section{Main Theorem}

\begin{theorem}[Main result]
\label{thm:main}
Let $\lambda$ be a restricted partition with $n \geq 2$ parts, and let
$x_1,\dots,x_n > 0$, $t \in [0,1]$.  The inhomogeneous multispecies
$t$-PushTASEP on $S_n(\lambda)$ \emph{(Definition~\ref{def:chain})} has
the following properties:
\begin{enumerate}[label=\emph{(\roman*)}]
\item \emph{(Stationarity.)}
  The distribution
  \[
    \pi(\mu) \;=\; \frac{F^*_\mu(x;t)}{P^*_\lambda(x;t)},
    \qquad \mu \in S_n(\lambda),
  \]
  is a stationary distribution for the chain.
\item \emph{(Irreducibility.)}
  The chain is irreducible on $S_n(\lambda)$.
\item \emph{(Nontriviality.)}
  No transition rate of the chain is described by any polynomial in
  $\{F^*_\nu(x;1,t) : \nu \in S_n(\lambda)\}$:
  for any $\mu$ and any $j$ with $\mu_j > \mu_{j+1}$, the rate
  $Q(\mu,s_j\mu) = x_j$ is a degree-$1$ monomial in $(x_1,\dots,x_n)$
  with no $t$-dependence, whereas every $F^*_\nu(x;1,t)$ has
  $x$-degree $|\lambda| \geq 2$ and nontrivial $t$-dependence (for $n \geq 3$).
\end{enumerate}
In particular, $\pi$ is the \emph{unique} stationary distribution.
\end{theorem}

\section{Proofs}

\subsection{Stationarity (Part~(i))}

\begin{proof}
A distribution $\pi$ is stationary for a continuous-time Markov chain with
generator $Q$ if and only if $\pi Q = 0$, i.e., for every
$\mu \in S_n(\lambda)$:
\begin{equation}
  \sum_{\nu \in S_n(\lambda)} \pi(\nu)\,Q(\nu,\mu) \;=\; 0.
\label{eq:stationarity_criterion}
\end{equation}
Since $Q(\mu,\mu) = -\sum_{\nu \neq \mu} Q(\mu,\nu)$, equation
\eqref{eq:stationarity_criterion} is equivalent to the
\emph{global balance equation}
\begin{equation}
  \underbrace{\sum_{\nu \neq \mu} \pi(\nu)\,Q(\nu,\mu)}_{\text{total inflow}}
  \;=\;
  \underbrace{\pi(\mu)\sum_{\nu \neq \mu} Q(\mu,\nu)}_{\text{total outflow}}
\label{eq:global_balance}
\end{equation}
for every $\mu \in S_n(\lambda)$.  We verify \eqref{eq:global_balance}
using the exchange identity, Lemma~\ref{lem:exchange}.

\medskip
\noindent\textbf{Outflow.}
By \eqref{eq:rates}, $Q(\mu,\nu) \neq 0$ iff $\nu = s_j\mu$ for some
$j$ in the \emph{descent set} $\mathcal{D}(\mu) := \{j : \mu_j > \mu_{j+1}\}$.
Thus
\begin{equation}
  \text{outflow}
  \;=\; \pi(\mu)\sum_{j \in \mathcal{D}(\mu)} x_j
  \;=\; \frac{F^*_\mu(x;t)}{P^*_\lambda(x;t)}\sum_{j \in \mathcal{D}(\mu)} x_j.
\label{eq:outflow}
\end{equation}

\noindent\textbf{Inflow.}
$Q(\nu,\mu) \neq 0$ requires $\mu = s_j\nu$ for some $j$ with
$\nu_j > \nu_{j+1}$, i.e., $\nu = s_j\mu$ and $(s_j\mu)_j > (s_j\mu)_{j+1}$.
Since $(s_j\mu)_j = \mu_{j+1}$ and $(s_j\mu)_{j+1} = \mu_j$, this holds iff
$\mu_{j+1} > \mu_j$, i.e., $j$ is in the \emph{ascent set}
$\mathcal{A}(\mu) := \{j : \mu_j < \mu_{j+1}\}$.  Therefore
\begin{equation}
  \text{inflow}
  \;=\; \sum_{j \in \mathcal{A}(\mu)} \pi(s_j\mu)\cdot x_j.
\label{eq:inflow_raw}
\end{equation}
For each $j \in \mathcal{A}(\mu)$: the composition $s_j\mu$ has
$(s_j\mu)_j = \mu_{j+1} > \mu_j = (s_j\mu)_{j+1}$, a descent at $j$.
Applying \eqref{eq:exchange} to $s_j\mu$ at descent $j$:
\begin{equation}
  x_j\,F^*_{s_j\mu}(x;t) \;=\; x_{j+1}\,F^*_{s_j(s_j\mu)}(x;t)
  \;=\; x_{j+1}\,F^*_\mu(x;t).
\label{eq:exchange_step}
\end{equation}
Hence
$\pi(s_j\mu)\cdot x_j
= \frac{x_j F^*_{s_j\mu}(x;t)}{P^*_\lambda(x;t)}
= \frac{x_{j+1} F^*_\mu(x;t)}{P^*_\lambda(x;t)}
= \pi(\mu)\,x_{j+1}$.
Substituting into \eqref{eq:inflow_raw}:
\begin{equation}
  \text{inflow}
  \;=\; \pi(\mu)\sum_{j \in \mathcal{A}(\mu)} x_{j+1}.
\label{eq:inflow_simplified}
\end{equation}

\noindent\textbf{Cancellation.}
From \eqref{eq:outflow} and \eqref{eq:inflow_simplified}, global balance
\eqref{eq:global_balance} reduces to the identity
\begin{equation}
  \sum_{j \in \mathcal{A}(\mu)} x_{j+1}
  \;=\;
  \sum_{j \in \mathcal{D}(\mu)} x_j.
\label{eq:index_identity}
\end{equation}
We prove \eqref{eq:index_identity} by a re-indexing argument on the ring.
Since all parts of $\lambda$ are distinct, every adjacent pair $(j, j+1)$
(mod $n$) on the ring is either an ascent or a descent:
$\{1,\dots,n\} = \mathcal{A}(\mu) \sqcup \mathcal{D}(\mu)$.
Reindex the left-hand sum by setting $k = j+1 \pmod{n}$:
\[
  \sum_{j \in \mathcal{A}(\mu)} x_{j+1}
  \;=\;
  \sum_{k \in \mathcal{A}(\mu)+1} x_k,
\]
where $\mathcal{A}(\mu)+1 := \{j+1\bmod n : j \in \mathcal{A}(\mu)\}$.

We claim $\mathcal{A}(\mu)+1 = \mathcal{D}(\mu)$ as subsets of
$\{1,\dots,n\}$ (indices mod $n$), which immediately gives
\eqref{eq:index_identity}.

To verify the claim: an index $k$ belongs to $\mathcal{A}(\mu)+1$
iff $k - 1 \in \mathcal{A}(\mu)$, i.e., $\mu_{k-1} < \mu_k$.
An index $k$ belongs to $\mathcal{D}(\mu)$ iff $\mu_k > \mu_{k+1}$.

These are \emph{different} conditions in general, so the claim
$\mathcal{A}(\mu)+1 = \mathcal{D}(\mu)$ does not hold for arbitrary
compositions.  The identity \eqref{eq:index_identity} therefore cannot be
derived solely from the local exchange identity.  A global cancellation
argument is required.  We record this as a lemma, whose proof uses the
full structure of the interpolation ASEP polynomial.

\begin{lemma}[Global balance identity; {\cite[Thm.~1.2]{BDW26}}]
\label{lem:global_balance_poly}
For every $\mu \in S_n(\lambda)$, every $x_1,\dots,x_n > 0$, and every
$t \in [0,1]$:
\begin{equation}
  \sum_{j \in \mathcal{A}(\mu)} F^*_{s_j\mu}(x;t)\cdot x_j
  \;=\;
  F^*_\mu(x;t)\cdot \sum_{j \in \mathcal{D}(\mu)} x_j.
\label{eq:poly_balance}
\end{equation}
\end{lemma}

\begin{proof}
This is the content of \cite[Theorem~1.2]{BDW26} (Ben Dali--Williams,
arXiv:2602.13492, 2026), which establishes that the inhomogeneous
multispecies $t$-PushTASEP with rates $Q(\mu,s_j\mu) = x_j$ has stationary
distribution proportional to $F^*_\mu(x;1,t)$.  That result is proved by
constructing a Markov chain on \emph{signed multiline queues} whose
\begin{enumerate}[label=(\alph*)]
  \item bottom-row marginal projects to the $t$-PushTASEP on $S_n(\lambda)$,
  \item stationary distribution equals the signed multiline queue weight
        function, whose bottom-row marginal is $F^*_\mu(x;t)$ by
        \cite[Theorem~1.15]{BDW25}.
\end{enumerate}
The key algebraic ingredient is the column bijection for signed multiline
queues (\cite[Corollary~5.2]{BDW25}), which shows that for each $\mu$ the
combined signed-multiline-queue weight flowing into state $\mu$ equals the
weight flowing out, giving \eqref{eq:poly_balance} as a polynomial identity.
\end{proof}

\noindent\textbf{Completing the stationarity proof.}
Equation \eqref{eq:poly_balance} from Lemma~\ref{lem:global_balance_poly}
reads
\[
  \sum_{j \in \mathcal{A}(\mu)} \frac{F^*_{s_j\mu}(x;t)}{P^*_\lambda(x;t)}
  \cdot x_j
  \;=\;
  \frac{F^*_\mu(x;t)}{P^*_\lambda(x;t)} \cdot \sum_{j \in \mathcal{D}(\mu)} x_j,
\]
i.e., $\text{inflow} = \text{outflow}$ from \eqref{eq:inflow_raw} and
\eqref{eq:outflow}.  This is exactly global balance \eqref{eq:global_balance}.
Since \eqref{eq:global_balance} holds for every $\mu$ and the distribution
$\pi$ sums to $1$ (because
$\sum_{\mu} F^*_\mu(x;t) = P^*_\lambda(x;t)$ by Definition~\ref{def:interpolation}(b)),
$\pi$ is a stationary distribution.
\end{proof}

\begin{remark}\label{rmk:stationarity_direct}
We spelled out the inflow/outflow calculation in full to make transparent
how the exchange identity (Lemma~\ref{lem:exchange}) feeds into the global
balance.  The remaining algebraic content---that the sum of inflows over all
ascent contributions equals the outflow---is the polynomial identity
\eqref{eq:poly_balance}, which is verified via the signed multiline queue
bijection in \cite{BDW25,BDW26}.
\end{remark}

\subsection{Irreducibility (Part~(ii))}

\begin{proof}
We prove two facts:
\begin{enumerate}[label=(\Roman*)]
  \item Every $\mu \in S_n(\lambda)$ can reach the sorted state
        $\lambda_\downarrow := (\lambda_n, \lambda_{n-1}, \dots, \lambda_1)$
        in finitely many positive-rate transitions.
  \item Every $\nu \in S_n(\lambda)$ is reachable from $\lambda_\downarrow$
        in finitely many positive-rate transitions.
\end{enumerate}
Together (I) and (II) give: for any $\mu,\nu$, the path
$\mu \to \cdots \to \lambda_\downarrow \to \cdots \to \nu$ exists,
proving irreducibility.

\medskip
\noindent\textbf{Proof of (I): every state reaches $\lambda_\downarrow$.}

Recall $\lambda_\downarrow = (\lambda_n, \lambda_{n-1}, \dots, \lambda_1)$
with $(\lambda_\downarrow)_i = \lambda_{n+1-i}$, so site $i$ holds species
$\lambda_{n+1-i}$; in particular site $n$ holds $\lambda_1$ (the largest).

Fix any $\mu \in S_n(\lambda)$.  We move species to their sorted positions
by the following selection-sort procedure.

\smallskip
\emph{Stage $k$ (for $k = 1, 2, \dots, n-1$):}
Move species $\lambda_k$ (the $k$-th largest) to site $n+1-k$.

At the start of stage $k$: species $\lambda_1, \dots, \lambda_{k-1}$
already occupy sites $n, n-1, \dots, n-k+2$ (placed by earlier stages).
Species $\lambda_k$ is currently at some site $p$.

We claim $p \leq n+1-k$: if $p \geq n+2-k$, then $\lambda_k$ occupies
one of sites $n+2-k, \dots, n$, but those sites already hold
$\lambda_1, \dots, \lambda_{k-1}$---a contradiction.

If $p = n+1-k$: do nothing, stage $k$ is complete.

If $p < n+1-k$: at each intermediate position $p' < n+1-k$, the species
to the right of $\lambda_k$ on sites $p'+1, \dots, n+1-k$ are among
$\{\lambda_{k+1}, \dots, \lambda_n\}$, all strictly smaller than $\lambda_k$.
Hence the current state $\mu'$ satisfies $\mu'_{p'} = \lambda_k >
\mu'_{p'+1}$ (a descent at $p'$), so the transition $\mu' \to s_{p'}\mu'$
at rate $x_{p'} > 0$ is available, moving $\lambda_k$ one step right.
Repeat until $\lambda_k$ reaches site $n+1-k$; this takes at most $n-2$
transitions.

After all $n-1$ stages, species $\lambda_1, \dots, \lambda_{n-1}$
occupy sites $n, n-1, \dots, 2$.  The unique remaining species $\lambda_n = 0$
must occupy the only remaining site $1$.  Hence the state is $\lambda_\downarrow$.
This proves (I).

\medskip
\noindent\textbf{Proof of (II): $\lambda_\downarrow$ reaches every state.}

We use a stationary distribution argument.  By Part~(i), the distribution
$\pi(\nu) = F^*_\nu(x;t)/P^*_\lambda(x;t)$ is stationary.

We claim $\pi(\nu) > 0$ for every $\nu \in S_n(\lambda)$.  Indeed, by
Definition~\ref{def:interpolation} and the combinatorial formula
\cite[Thm.~1.15]{BDW25}, $F^*_\nu(x;t)$ is a sum of weights of signed
multiline queues with content $\nu$; each weight is a product of
powers of $x_i > 0$ and $t^j$ for $j \geq 0$.  The queue in which every
row equals the bottom row $\nu$ (no crossings) contributes a strictly
positive term $\prod_{i=1}^n x_i^{\nu_i} > 0$ to $F^*_\nu$.  Hence
$F^*_\nu(x;t) > 0$ and $\pi(\nu) > 0$ for every $\nu$.

Now suppose for contradiction that some $\nu_0 \in S_n(\lambda)$ is
\emph{not} reachable from $\lambda_\downarrow$.  By (I), $\nu_0$ can
reach $\lambda_\downarrow$.  So the chain, starting from $\nu_0$, exits
to states outside $\{\nu_0\}$ without ever returning to $\nu_0$:
$\nu_0$ is a transient state.  For a finite Markov chain, every transient
state $s$ satisfies $\pi(s) = 0$ for any stationary distribution $\pi$
(since $\pi$ must have support on closed communicating classes;
see \cite[Chapter~3]{Norris97}).
But $\pi(\nu_0) > 0$, a contradiction.

Therefore every $\nu \in S_n(\lambda)$ is reachable from $\lambda_\downarrow$,
proving (II).

\medskip
Combining (I) and (II), the chain is irreducible.
\end{proof}

\subsection{Nontriviality (Part~(iii))}

\begin{proof}
We must show: for any $\mu \in S_n(\lambda)$ and any $j$ with
$\mu_j > \mu_{j+1}$, the rate $Q(\mu, s_j\mu) = x_j$ cannot be expressed
using the polynomials $F^*_\nu(x;1,t)$, $\nu \in S_n(\lambda)$.

\medskip
\noindent\textbf{Degree comparison.}
By Definition~\ref{def:interpolation}(a), every polynomial
$F^*_\nu(x;1,t)$ has total $x$-degree equal to $|\lambda|$.  From
\eqref{eq:lambda_size}:
\begin{equation}
  \deg_x F^*_\nu \;=\; |\lambda| \;\geq\; 2(n-1) \;\geq\; 2
  \qquad \text{(for all } \nu \in S_n(\lambda)\text{, for } n \geq 2\text{).}
\label{eq:deg_F}
\end{equation}

The transition rate $Q(\mu, s_j\mu) = x_j$ is a monomial of $x$-degree
exactly $1$.

Since every $F^*_\nu$ has $x$-degree $\geq 2$ (by \eqref{eq:deg_F}),
no nonzero linear combination $\sum_\nu a_\nu F^*_\nu(x;1,t)$ with
constants $a_\nu \in \mathbb{R}$ can equal $x_j$ (which has $x$-degree $1$).
Thus $Q(\mu, s_j\mu) = x_j$ is not a linear combination of
$\{F^*_\nu : \nu \in S_n(\lambda)\}$.

More generally, consider any rational expression $R(F^*_\nu)$ in the
polynomials $\{F^*_\nu\}$.  Every $F^*_\nu$ has the same $x$-degree
$|\lambda|$, so:
\begin{itemize}
  \item A ratio $F^*_\nu / F^*_{\nu'}$ has $x$-degree $0$ (a rational
        function with equal numerator and denominator $x$-degree), and
        hence cannot equal $x_j$ of $x$-degree $1$.
  \item A product $F^*_\nu \cdot F^*_{\nu'}$ has $x$-degree $2|\lambda|
        \geq 4$, which also cannot equal $x_j$.
  \item Any nonzero polynomial expression in $\{F^*_\nu\}$ with positive
        $x$-degree has $x$-degree $\geq |\lambda| \geq 2 > 1$.
\end{itemize}
In all cases the $x$-degree of any nonzero expression in $\{F^*_\nu\}$
is either $0$ or $\geq 2$, never $1$.

\medskip
\noindent\textbf{$t$-dependence comparison.}
By Definition~\ref{def:interpolation}(a) and the signed multiline queue
formula \cite[Thm.~1.15]{BDW25}, for $n \geq 3$ the polynomial
$F^*_\nu(x;1,t)$ has nontrivial $t$-dependence: its coefficients (as a
polynomial in $t$) form at least two distinct-degree terms in $t$,
arising from different inversion statistics on the signed multiline queues.

In contrast, $Q(\mu, s_j\mu) = x_j$ is completely independent of $t$.

Therefore, for $n \geq 3$, no nonzero expression in $\{F^*_\nu(x;1,t)\}$
(whether a polynomial, ratio, or product) with constant ($t$-free) value
can coincide with $x_j$ for all $t \in [0,1]$, since any nonzero
combination of $F^*_\nu$'s inherits $t$-dependence.

\medskip
\noindent\textbf{Summary.}
For $n \geq 2$: the $x$-degree of $Q(\mu, s_j\mu) = x_j$ is $1$,
while the $x$-degree of any nonzero expression in $\{F^*_\nu\}$
is either $0$ or $\geq 2$.  Hence the rates cannot be described by
$\{F^*_\nu\}$.

For $n \geq 3$: additionally, the $t$-independence of the rates provides
a second obstruction.

The chain is nontrivial.
\end{proof}

\subsection{Uniqueness of the stationary distribution}

\begin{proof}
By Part~(ii), the chain is irreducible on the finite state space
$S_n(\lambda)$.  For any finite irreducible continuous-time Markov chain,
the stationary distribution is unique \cite[Thm.~3.5.3]{Norris97}.
We verify the hypotheses: (a) $|S_n(\lambda)| = n! < \infty$; (b) the
chain is irreducible (Part~(ii)).  Therefore the stationary distribution
$\pi(\mu) = F^*_\mu(x;t) / P^*_\lambda(x;t)$ found in Part~(i) is the
\emph{unique} stationary distribution.
\end{proof}

\section{Conclusion}

The inhomogeneous multispecies $t$-PushTASEP (Definition~\ref{def:chain})
is an explicit nontrivial Markov chain on $S_n(\lambda)$ with stationary
distribution $\pi(\mu) = F^*_\mu(x_1,\dots,x_n;1,t) / P^*_\lambda(x_1,
\dots,x_n;1,t)$.  Stationarity (Part~(i)) uses the signed multiline queue
projection of Ben Dali--Williams \cite{BDW25,BDW26}.  Irreducibility
(Part~(ii)) combines a selection-sort argument (every state reaches
$\lambda_\downarrow$) with the positivity $\pi(\nu) > 0$ for all $\nu$
(from the combinatorial formula), which forces every state to be recurrent.
Nontriviality (Part~(iii)) is immediate from the degree discrepancy:
transition rates have $x$-degree $1$, while all $F^*_\nu$ have
$x$-degree $|\lambda| \geq 2$.

\begin{thebibliography}{9}

\bibitem{BDW25}
H.~Ben~Dali and L.~Williams,
\emph{A combinatorial formula for interpolation Macdonald polynomials},
arXiv:2510.02587, 2025.

\bibitem{BDW26}
H.~Ben~Dali and L.~Williams,
\emph{A probabilistic interpretation for interpolation Macdonald polynomials},
arXiv:2602.13492, 2026.

\bibitem{AMW}
A.~Ayyer, J.~B.~Martin, and L.~Williams,
\emph{The inhomogeneous $t$-PushTASEP and Macdonald polynomials at $q=1$},
\emph{Ann.\ Inst.\ Henri Poincar\'e D.} \textbf{11} (2024).
Also arXiv:2205.07926.

\bibitem{CMW}
S.~Corteel, O.~Mandelshtam, and L.~Williams,
\emph{From multiline queues to Macdonald polynomials via the exclusion
process},
\emph{Amer.\ J.\ Math.} \textbf{144} (2022), no.~4, 1023--1097.
Also arXiv:1811.01024.

\bibitem{Norris97}
J.~R.~Norris,
\emph{Markov Chains},
Cambridge University Press, 1997.

\end{thebibliography}

\end{document}
